\subsection{全序集}

定义9. 两个元素 \(x, y\) （预序集 \(E\) 的元素）称为可比较的，如果关系“ \(x \leqslant y\) 或 \(y \leqslant x\) ”为真。一个集合 \(E\) 称为全序的，如果它是有序的，并且 \(E\) 的任意两个元素是可比较的。 \(E\) 上的序则被称为全序，相应的序关系称为全序关系。

如果 \(x\) 与 \(y\) 为全序集的元素，则要么 \(x = y\) ，要么 \(x < y\) ，要么 \(x > y\) ； \(x \leqslant y\) 的否定因而是 \(x > y\) .

 \(\mathbf{E}\) 上的一个序关系是全序的，当且仅当其图 \(\mathbf{G}\) 满足关系 \(\mathbf{G} \cup \overline{\mathbf{G}} = \mathbf{E} \times \mathbf{E}\) 以及关系 \(\mathbf{G} \circ \mathbf{G} = \mathbf{G}\) 和 \(\mathbf{G} \cap \overline{\mathbf{G}} = \Delta\) .

示例

\begin{enumerate}
\def\labelenumi{(\arabic{enumi})}
\item
  全序集的每个子集在诱导序下都是全序的。
\item
  设 E 为任意有序集。E 的空子集是全序的，E 的每个单元素子集也是全序的。
\end{enumerate}

* (3) 集合 \(\mathbf{R}\) （实数集）是全序的。*

\begin{enumerate}
\def\labelenumi{(\arabic{enumi})}
\setcounter{enumi}{3}
\tightlist
\item
  若 A 是至少含有两个不同元素的集合，则集合 \(\mathfrak{P}(A)\) （按包含关系排序）不是全序的，因为若 \(x \neq y\) ，则子集 \(\{x\}\) 与 \(\{y\}\) 不可比较。
\end{enumerate}

全序集在相反序下也是全序的；它是一个格，因而更是左有向且右有向的。

命题 11. 每个严格单调映射 \(f\) （从全序集 \(\mathbf{E}\) 到有序集 \(\mathbf{F}\) 的映射）都是单射。若 \(f\) 严格递增，则 \(f\) 是 \(\mathbf{E}\) 到 \(f(\mathbf{E})\) 上的一个同构.

因为 \(x \neq y\) 蕴含 \(x < y\) 或 \(x > y\) ；因此

\[
f (x) <   f (y) \quad \text { 或 } \quad f (x) > f (y),
\]

使得 \(f(x) \neq f(y)\) 在任一情形下均如此。还需证明若 \(f\) 严格递增，则 \(f(x) \leqslant f(y)\) 蕴含 \(x \leqslant y\) ；否则，我们应有 \(x > y\) ，从而 \(f(x) > f(y)\) .

命题 12. 设 E 为全序集，X 为 E 的子集。对于 E 中的元素 b 而言，b 是 X 在 E 中的最小上界，当且仅当 (1) b 是 X 的上界，(2) 对每个 \(c \in E\) （使得 c \textless{} b ），存在 \(x \in X\) 使得 \(c < x \leqslant b\) .

第二个条件说明没有元素 c \textless{} b 是 X 的上界，即 b 是 X 的上界集合 M 的极小元素；但由于 M 是全序的（第 10 号，命题 10），这与 b 是 M 的最小元素是等价的。

